\section{双元、单元、高次关系与隐函数积分全流程}\label{ic:indefinite_dual_complete}

\begin{examtip}本节把双元名称还原为约束微分、齐次化与有理参数。有理积分采用记号$\mathcal U$，各系数由微分推导。\end{examtip}

\subsection{圆型、双曲型与齐次比值}\label{ic:indefinite_dual_complete:differentials}\begin{method}{圆型、双曲型与齐次比值}\textbf{识别信号} 两个可微量$u,v$满足$u^2+v^2=a^2$或$v^2-u^2=a^2$。\par \textbf{操作步骤} 圆型取角或比值；双曲型取$u+v$；同时记录约束的微分。商导数给$\dd(\frac{u}{v})$，乘积导数给分部闭合。\par \textbf{疑难点与条件} 角变量公式局部成立，跨象限须补连续分支；根式符号不能任意取消。\end{method}

\begin{theorem}{角微分与双曲对数微分}设$a>0$，并在$v\ne0$的可微实支上讨论。若$u^2+v^2=a^2$，则$\dd\arctan(\frac{u}{v})=\frac{v\dd u-u\dd v}{a^2}=\frac{\dd u}{v}$，在$v\ne0$的局部角支成立。若$v^2-u^2=a^2$，则$\dd\ln|u+v|=\frac{\dd u}{v}$。\end{theorem}\begin{proof}分别用反正切导数和约束$v\dd v=\mp u\dd u$计算。圆型中$v\dd u-u\dd v=\frac{(v^2+u^2)\dd u}{v}$；双曲型$\frac{\dd u+\dd v}{u+v}=\frac{\dd u}{v}$。\end{proof}

\begin{example}{考研风格｜圆型非齐次分式}{难度：$\star\star\star$}\label{ic:ex:indefinite_dual_complete:circlefraction}\textbf{题目} 求$\int\frac{\dd x}{\sqrt{1-x^2}(1+\sqrt{1-x^2})}$，$|x|<1$。\par \textbf{思路} 先取圆的角参数，再半角。\par \textbf{详细解答} $x=\sin t$，$t\in(-\frac{\pi}{2},\frac{\pi}{2})$，积分$\int\frac{\dd t}{1+\cos t}=\tan(\frac{t}{2})+C=\frac{x}{1+\sqrt{1-x^2}}+C$。另一法直接对商求导，利用$x^2+(1-x^2)=1$整理。\par \textbf{求导验算} 若$v=\sqrt{1-x^2}$，$[\frac{x}{1+v}]'=\frac{1+v+\frac{x^2}{v}}{(1+v)^2}=\frac{1}{v(1+v)}$。\par \textbf{疑难点与方法比较} 先除去一个圆型微分，再处理常数加双元，比分母全体有理化更短。\end{example}

\begin{example}{竞赛风格｜双元原函数再分部}{难度：$\star\star\star$}\label{ic:ex:indefinite_dual_complete:dual_log}\textbf{题目} 求$\int\ln(x+\sqrt{1+x^2})\dd x$。\par \textbf{思路} 对数的导数已由双元公式确定。\par \textbf{详细解答} 令$v=\sqrt{1+x^2}$，分部得$x\ln(x+v)-\int \frac{x\dd x}{v}=x\ln(x+v)-v+C$。也可$x=\sinh t$，积分$\int t\cosh t\dd t=t\sinh t-\cosh t+C$。\par \textbf{求导验算} 导数为$\ln(x+v)+\frac{x}{v}-\frac{x}{v}$。\par \textbf{疑难点与方法比较} $x+v>0$，此处实对数无需绝对值。\end{example}

\subsection{恒等式诱导与莫比乌斯变换的构造}\label{ic:indefinite_dual_complete:mobius}\begin{method}{恒等式诱导与莫比乌斯变换的构造}\textbf{识别信号} 分母含一次式与二次根式的乘积。\par \textbf{操作步骤} 把二次式写成两个一次式平方和或差，分别除以同一个一次式平方；所产生的两个比值满足圆或双曲关系；计算比值微分。\par \textbf{疑难点与条件} 不能只写“注意到恒等式”，须给平方展开或系数比较。\end{method}

\begin{example}{竞赛风格｜一次分母乘二次根式双解}{难度：$\star\star\star$}\label{ic:ex:indefinite_dual_complete:mobiusroot}\textbf{题目} 求$\int\frac{\dd x}{(x+1)\sqrt{x^2+1}}$，$x\ne-1$。\par \textbf{思路} 用$2(x^2+1)=(x+1)^2+(x-1)^2$产生双曲双元。\par \textbf{详细解答} 设$u=\frac{x-1}{x+1},v=\frac{\sqrt2\sqrt{x^2+1}}{x+1}$，有$v^2-u^2=1$，且$\frac{\dd u}{v}=\frac{\sqrt2\,\dd x}{(x+1)\sqrt{x^2+1}}$。所以原函数$\frac1{\sqrt2}\ln\left|\frac{x-1+\sqrt2\sqrt{x^2+1}}{x+1}\right|+C$。第二法$x=\sinh t$，积分$\int\frac{\dd t}{1+\sinh t}$，再$\e^t$化有理式，可得等价对数。\par \textbf{求导验算} 由$\dd\ln|u+v|=\frac{\dd u}{v}$直接验证，$v$可正可负但不为0。\par \textbf{疑难点与方法比较} 以带符号的$v$保留分支，比无条件替成正平方根安全。\end{example}

\begin{example}{竞赛风格｜四次根式的平方分解}{难度：$\star\star\star$}\label{ic:ex:indefinite_dual_complete:quarticidentity}\textbf{题目} 求$\int\frac{\dd x}{x\sqrt{1+x^2+3x^4}}$，$x>0$。\par \textbf{思路} 除以$x^4$后配方，或先倒数平方换元。\par \textbf{详细解答} 令$t=x^{-2}$，$\dd t=-2x^{-3}\dd x$，原式$=-\frac12\int\frac{\dd t}{\sqrt{t^2+t+3}}$。配方得$-\frac12\ln[t+\frac{1}{2}+\sqrt{t^2+t+3}]+C$。等价地利用$4(1+x^2+3x^4)=(x^2+2)^2+11x^4$构造双曲双元，得到相同对数形式。\par \textbf{求导验算} 对$t$的原函数导数为$-\frac{1}{2\sqrt{t^2+t+3}}$，乘$t'$恢复被积式。\par \textbf{疑难点与方法比较} 分母的$x$是倒数平方换元的信号；去掉它则不能沿用本解。\end{example}

\subsection{对勾三元与回文多项式}\label{ic:indefinite_dual_complete:reciprocal}\begin{method}{对勾三元与回文多项式}\textbf{识别信号} 同时出现$x^2$与$x^{-2}$，或四次回文多项式。\par \textbf{操作步骤} 设$p=x+\frac{a}{x},q=x-\frac{a}{x}$，则$p^2-q^2=4a$；根据微分中的$1\mp \frac{a}{x^2}$选变量；根式除以$|x|$后识别二次式。\par \textbf{疑难点与条件} 回文结构不保证任意分子都能直接换元；原微分的匹配决定是否初等化。\end{method}

\begin{example}{竞赛风格｜同一四次根式的两种信号}{难度：$\star\star\star$}\label{ic:ex:indefinite_dual_complete:reciprocalplus}\textbf{题目} 求$\int\frac{(x^2+1)\dd x}{x\sqrt{x^4+5x^2+1}}$，$x>0$。\par \textbf{思路} 选$q=x-\frac{1}{x}$，因为$q'=1+x^{-2}$。\par \textbf{详细解答} $\sqrt{x^4+5x^2+1}=x\sqrt{q^2+7}$，故积分$\int\frac{\dd q}{\sqrt{q^2+7}}=\ln(q+\sqrt{q^2+7})+C$。如果分子改为$x^2-1$，改取$p=x+\frac{1}{x}$，得到$\ln(p+\sqrt{p^2+3})+C$。\par \textbf{求导验算} 分别用$q'=1+x^{-2},p'=1-x^{-2}$验证两题。\par \textbf{疑难点与方法比较} 分子正负决定用和或差；不能凭分母回文就固定一种换元。\end{example}

\begin{example}{竞赛风格｜对勾结构导出反正切}{难度：$\star\star\star$}\label{ic:ex:indefinite_dual_complete:reciprocalangle}\textbf{题目} 求$\int\frac{(1-x^2)\dd x}{(1+x^2)\sqrt{x^4+x^2+1}}$。\par \textbf{思路} 设$r=\sqrt{x^4+x^2+1}$，试角变量$\frac{x}{r}$。\par \textbf{详细解答} 计算$(\frac{x}{r})'=\frac{1-x^4}{r^3}$，而$1+(\frac{x}{r})^2=\frac{(x^2+1)^2}{r^2}$。所以被积函数等于$[\arctan(\frac{x}{r})]'$，答案$\arctan(\frac{x}{\sqrt{x^4+x^2+1}})+C$。也可先以$x\pm\frac{1}{x}$构造三元，约成$\int\frac{\dd u}{1+u^2}$。\par \textbf{求导验算} 上面两条导数比值化简为$\frac{1-x^2}{(1+x^2)r}$，在$x=0$也可直接连续验证。\par \textbf{疑难点与方法比较} 把直接验证看作第二层证据，不能省略识别出角变量的恒等式。\end{example}

\subsection{双元化单元与根式和差参数}\label{ic:indefinite_dual_complete:single}\begin{method}{双元化单元与根式和差参数}\textbf{识别信号} 双元和差乘积为常数，或多个有理分数幂共用基底。\par \textbf{操作步骤} 用和$t=u+v$并由乘积反解差；嵌套根式再取合适根次；多个分数幂取分母最小公倍数。\par \textbf{疑难点与条件} 参数须选与原根式一致的实支；幂换元对负自变量需另检。\end{method}

\begin{example}{竞赛风格｜根式和差同时参数化}{难度：$\star\star\star$}\label{ic:ex:indefinite_dual_complete:nestedtwo}\textbf{题目} 求$\int[\sqrt{x+1}-\sqrt x]^{-2}\dd x$，$x>0$。\par \textbf{思路} 令$t=\sqrt{x+1}+\sqrt x$，其倒数是两根之差。\par \textbf{详细解答} $\sqrt{x+1}=\frac{t+t^{-1}}{2},\sqrt x=\frac{t-t^{-1}}{2}$，所以$x=\frac{t^2-2+t^{-2}}{4}$、$\dd x=\frac{(t-t^{-3})\dd t}{2}$。原式$\frac12\int(t^3-t^{-1})\dd t=\frac{t^4}{8}-\frac12\ln t+C$。另一法先将被积式有理化为$(\sqrt{x+1}+\sqrt x)^2$，再分别积分，结果相差常数。\par \textbf{求导验算} 乘参数导数可回到$t^2\dd x$，而$t^2$正是原被积函数。\par \textbf{疑难点与方法比较} 同一和差参数同时消掉两根，避免逐层换元。\end{example}

\subsection{高次双元：统一微分与有理参数}\label{ic:indefinite_dual_complete:highdual}\begin{method}{高次双元：统一微分与有理参数}\textbf{识别信号} $y^n=\varepsilon x^n+c$，$n\ge2,\varepsilon=\pm1,c\ne0$。\par \textbf{操作步骤} 取$t=\frac{x}{y}$，用约束微分计算$\dd t$；把$y^n$消成$\frac{c}{1-\varepsilon t^n}$；所需积分归到有理函数族$\mathcal U$。\par \textbf{疑难点与条件} $c=0$退化为幂函数关系；不能代入含$\frac{1}{c}$的公式。\end{method}

\begin{theorem}{高次双元的基本式与降次式}在可微实支上，$\dd(\frac{x}{y})=\frac{c\,\dd x}{y^{n+1}}$，$\int\frac{\dd x}{y}=\int\frac{\dd t}{1-\varepsilon t^n}$。另记$J_m=\int y^m\dd x$，$m\ne-1$时$J_m=\frac{xy^m}{m+1}+\frac{mc\,J_{m-n}}{m+1}$。\end{theorem}\begin{proof}由$y'=\frac{\varepsilon x^{n-1}}{y^{n-1}}$得商导数$\frac{c}{y^{n+1}}$，且$y^n(1-\varepsilon t^n)=c$。再算$(xy^m)'=y^m+m\varepsilon x^ny^{m-n}=(m+1)y^m-mc\,y^{m-n}$，移项得到递推。\end{proof}

\begin{example}{竞赛风格｜四次双元的初等化}{难度：$\star\star\star$}\label{ic:ex:indefinite_dual_complete:fourthdual}\textbf{题目} 求$\int\frac{\dd x}{(1+x^4)^{\frac{5}{4}}}$及$\int\frac{\dd x}{(1+x^4)^{\frac{1}{4}}}$。\par \textbf{思路} 两者用同一$t=\frac{x}{(1+x^4)^{\frac{1}{4}}}$，但微分系数不同。\par \textbf{详细解答} 取$y=(1+x^4)^{\frac{1}{4}}>0$，第一题由$\dd t=\frac{\dd x}{y^5}$直接得$t+C$。第二题化为$\int\frac{\dd t}{1-t^4}$；部分分式为$\frac12(1-t^2)^{-1}+\frac12(1+t^2)^{-1}$，答案$\frac14\ln|\frac{1+t}{1-t}|+\frac12\arctan t+C$。此处$|t|<1$。\par \textbf{求导验算} 第二式对$t$导数为$\frac{1}{1-t^4}$，再乘$\frac{1}{y^5}$，因$1-t^4=\frac{1}{y^4}$得$\frac{1}{y}$。\par \textbf{疑难点与方法比较} 高次根式也可初等；关键是指数与约束的匹配。\end{example}

\begin{example}{竞赛风格｜高次双元递推的系数}{难度：$\star\star\star$}\label{ic:ex:indefinite_dual_complete:cubicrecur}\textbf{题目} 求$\int(x^3+2)^{\frac{2}{3}}\dd x$，取$x>0$。\par \textbf{思路} $n=3,c=2,m=2$，降到$J_{-1}$。\par \textbf{详细解答} 令$y=(x^3+2)^{\frac{1}{3}},t=\frac{x}{y}$。递推给$J_2=\frac{xy^2}{3}+\frac{4J_{-1}}{3}$。而$J_{-1}=\int\frac{\dd t}{1-t^3}=-\frac13\ln(1-t)+\frac16\ln(t^2+t+1)+\frac1{\sqrt3}\arctan(\frac{2t+1}{\sqrt3})$。代入即为完整原函数。\par \textbf{求导验算} 由已证递推和$t'=\frac{2}{y^4}$核验；$1-t^3=\frac{2}{y^3}$消去全部余项。\par \textbf{疑难点与方法比较} 递推中的$mc$不能漏掉；各项系数应由求导核验，避免幂次和下标混淆。\end{example}

\subsection{统一有理积分族与公因数降阶}\label{ic:indefinite_dual_complete:ufamily}\begin{method}{统一有理积分族与公因数降阶}\textbf{识别信号} $\mathcal U^\varepsilon_{m,n}(x)=\int \frac{x^{m-1}}{1+\varepsilon x^n}\dd x$，$m,n$为正整数。\par \textbf{操作步骤} $m\ge n$先长除；有公因数先$x^k$降阶；剩余真分式按单位根配共轭，转成实对数及反正切。\par \textbf{疑难点与条件} 降阶有$\frac{1}{k}$系数；该记号表示可计算有理积分，不是新的非初等函数。\end{method}

\begin{theorem}{公因数降阶的比例因子}$\mathcal U^\varepsilon_{km,kn}(x)=k^{-1}\mathcal U^\varepsilon_{m,n}(x^k)+C$。若$P=1+\varepsilon x^n$有互异复根$r_j$，则$\frac{x^{m-1}}{P}=\sum \frac{r_j^{m-1}}{P'(r_j)(x-r_j)}$（$m\le n$）。\end{theorem}\begin{proof}第一式令$u=x^k$，$x^{km-1}\dd x=\frac{u^{m-1}\dd u}{k}$。第二式由各简单极点的系数及真分式唯一性得到；实系数下共轭项合并为实表达。\end{proof}

\begin{example}{考研风格｜公因数与长除联用}{难度：$\star\star$}\label{ic:ex:indefinite_dual_complete:gcd}\textbf{题目} 求$\int \frac{x^5\dd x}{1+x^4}$。\par \textbf{思路} 先$x^2=t$，再除法。\par \textbf{详细解答} $x^5\dd x=\frac{t^2\dd t}{2}$，积分$\frac12\int[1-\frac{1}{1+t^2}]\dd t=\frac{x^2}2-\frac12\arctan(x^2)+C$。也可$\frac{x^5}{1+x^4}=x-\frac{x}{1+x^4}$后用$x^2$凑微分。\par \textbf{求导验算} 导数$x-\frac{x}{1+x^4}=\frac{x^5}{1+x^4}$。\par \textbf{疑难点与方法比较} 公因数降阶必须同时变换分子幂和微分。\end{example}

\subsection{隐函数：参数化与约束微分}\label{ic:indefinite_dual_complete:implicit}\begin{method}{隐函数：参数化与约束微分}\textbf{识别信号} $F(x,y)=0$且所取支路可微。\par \textbf{操作步骤} 先找$t=x-y,\frac{y}{x}$等使曲线可有理参数化的组合；或由$\dd F=0$把原微分改成$\dd y$，利用分部和原约束消元。\par \textbf{疑难点与条件} 若$F_y=0$，不能自动视为可微$y(x)$；原被积函数的分母零点也须排除。\end{method}

\begin{example}{竞赛风格｜同一隐函数积分的两解}{难度：$\star\star\star$}\label{ic:ex:indefinite_dual_complete:implicitcurve}\textbf{题目} 若$y(x-y)^2=x$，求$\int\frac{\dd x}{x-3y}$，在正则且分母非零的支路上。\par \textbf{思路} 参数$t=x-y$能同时解出$x,y$。\par \textbf{详细解答} 由$y(t^2-1)=t$得$y=\frac{t}{t^2-1},x=\frac{t^3}{t^2-1}$，因此$\frac{\dd x}{x-3y}=\frac{t\dd t}{t^2-1}$，原函数$\frac12\ln|(x-y)^2-1|+C$。第二式由约束可得$\frac14[\ln|x|-3\ln|y|]+C$，因$\frac{x}{y^3}=(t^2-1)^2$与前式相同。\par \textbf{求导验算} 参数导数$x'=\frac{t^2(t^2-3)}{(t^2-1)^2}$与$x-3y=\frac{t(t^2-3)}{t^2-1}$约分验证。\par \textbf{疑难点与方法比较} 参数式的$t=0,\pm1,\pm\sqrt3$需按原曲线与被积式逐一检查，不能凭约分把奇点补成同一道题。\end{example}

\begin{example}{竞赛风格｜隐式三次关系的两种表示}{难度：$\star\star\star$}\label{ic:ex:indefinite_dual_complete:implicitcubic}\textbf{题目} $z^3+xz=8$定义可微支$z(x)$，求$\int z^2\dd x$。\par \textbf{思路} 把$x$用$z$表示，或构造含$xz^2$的分部关系。\par \textbf{详细解答} 因$z\ne0$，$x=\frac{8}{z}-z^2$，$\dd x=(-\frac{8}{z^2}-2z)\dd z$，故积分$=-8z-\frac{z^4}{2}+C$。等价形式$-xz^2-\frac32z^4+C$由原约束相乘得到。\par \textbf{求导验算} 由$(3z^2+x)z'=-z$验证第一式导数为$z^2$；只在$3z^2+x\ne0$的正则段使用。\par \textbf{疑难点与方法比较} 不要把隐式变量当常数；先变参数可省去复杂根式解。\end{example}
